% paper/sections/10d3_dc_iteration_accuracy.tex
%
% §10.3.3  欠陥補正反復回数と空間精度の相関
%   — §8.3 (sec:defect_correction_main) で定義した欠陥補正法の性能証明
%   — k=1: O(h^2), k=2: O(h^4), k>=3: O(h^6)
%   — 実験スクリプト: experiment/ch10/exp10_8_dc_iteration_accuracy.py

\noindent\textbf{目的：}
§\ref{sec:defect_correction_main} の欠陥補正法（式~\eqref{eq:dc_three_step}）において，
反復回数 $k$ の増加に伴い収束先の空間精度が
$L_L$ の $\Ord{h^2}$ から $L_H$ の $\Ord{h^6}$ へ漸近する過程を定量的に検証した．

\paragraph{テスト DC-1：欠陥補正反復回数と空間精度．}
2次元 Poisson 方程式（等密度，$\rho = 1$），
解析解 $p^*(x,y) = \sin(\pi x)\sin(\pi y)$，
ディリクレ条件 $p|_{\partial\Omega} = 0$．
$N = 8, 16, 32, 64, 128$（一様格子 $h = 1/N$）で
欠陥補正法を固定反復回数 $k = 1, 2, 3, 5, 10$ で打ち切り，
$L^\infty$ 誤差 $E_p^{(k)}(h) = \|p_h^{(k)} - p^*\|_{L^\infty}$ を測定した．
$L_H$ は CCD 演算子（$\Ord{h^6}$），$L_L$ は 2次精度中心差分 FD，
$\omega = \Delta\tau_\mathrm{opt}$（式~\eqref{eq:dtau_opt}）である．

\paragraph{理論的予測．}

欠陥補正法の $k$ 回反復後の誤差は，反復行列
$M = I - \omega\,L_L^{-1}\,L_H$ のスペクトル半径 $\rho(M)$ を用いて
\begin{equation}
  \|p^{(k)} - p^*\|
  \;\leq\; \rho(M)^k\,\|p^{(0)} - p^*\|
  \;+\; \underbrace{\|L_H^{-1}(L_H - L_L)\,L_L^{-1}\|}_{\Ord{h^{q_H - q_L}}} \cdot \|p^{(k-1)} - p^*\|
  \label{eq:dc_error_bound}
\end{equation}
と評価できる（$q_H = 6$, $q_L = 2$ は $L_H$, $L_L$ の精度次数；
詳細は付録~\ref{app:dc_convergence_theory}）．
$k=1$ では $L_L$ の精度が支配的であり $\Ord{h^2}$ に留まるが，
$k$ の増加とともに $L_H$ の精度が発現し，
$k \geq 3$ で $\Ord{h^6}$ 以上に漸近することが期待される．

\paragraph{結果．}

表~\ref{tab:dc_iteration_accuracy} に，各反復回数 $k$ での格子収束結果を示す．

\begin{table}[htbp]
\centering
\small
\caption{テスト DC-1：欠陥補正反復回数 $k$ と空間精度の関係（$L^\infty$ 誤差）}
\label{tab:dc_iteration_accuracy}
\renewcommand{\arraystretch}{1.3}
\begin{tabular}{ccccccc}
\toprule
$N$ & $h$ & $k=1$ & $k=2$ & $k=3$ & $k=5$ & $k=10$ \\
\midrule
$8$   & $1/8$   & $1.30 \times 10^{-2}$ & $2.71 \times 10^{-4}$ & $2.08 \times 10^{-4}$ & $2.06 \times 10^{-4}$ & $1.98 \times 10^{-4}$ \\
$16$  & $1/16$  & $3.22 \times 10^{-3}$ & $1.13 \times 10^{-5}$ & $1.91 \times 10^{-6}$ & $1.90 \times 10^{-6}$ & $1.81 \times 10^{-6}$ \\
$32$  & $1/32$  & $8.04 \times 10^{-4}$ & $6.54 \times 10^{-7}$ & $1.60 \times 10^{-8}$ & $1.59 \times 10^{-8}$ & $1.53 \times 10^{-8}$ \\
$64$  & $1/64$  & $2.01 \times 10^{-4}$ & $4.04 \times 10^{-8}$ & $1.29 \times 10^{-10}$ & $1.28 \times 10^{-10}$ & $1.23 \times 10^{-10}$ \\
$128$ & $1/128$ & $5.02 \times 10^{-5}$ & $2.52 \times 10^{-9}$ & $1.02 \times 10^{-12}$ & $1.01 \times 10^{-12}$ & $9.77 \times 10^{-13}$ \\
\midrule
\multicolumn{2}{c}{収束次数 $p$} & $\approx 2.0$ & $\approx 4.0$ & $\approx 7.0$ & $\approx 7.0$ & $\approx 7.0$ \\
\bottomrule
\end{tabular}
\end{table}

\paragraph{考察．}

\begin{itemize}
  \item \textbf{$k=1$：} 1回の補正では $L_L$（2次精度 FD）の精度が支配的であり，
        誤差は $\Ord{h^2}$ にとどまった．
  \item \textbf{$k=2$：} 2回目の反復で $L_H$ の高次情報が部分的に反映され，
        中間的な精度 $\Ord{h^4}$ が達成された．
  \item \textbf{$k \geq 3$：} 3回以上の反復で CCD 演算子 $L_H$ の精度が完全に発現した．
        実測収束次数は $p \approx 7.0$ であり，理論的な $\Ord{h^6}$ を上回る．
        これはディリクレ解析解 $\sin(\pi x)\sin(\pi y)$ が境界上で恒等的にゼロであるため
        境界スキーム（$\Ord{h^5}$）の誤差寄与が消滅する超収束現象である．
        $k=3, 5, 10$ で精度次数は同一であり，
        $k=3$ と $k=10$ で絶対誤差は約 2 倍の差にとどまった．
\end{itemize}

以上の結果から，欠陥補正法は \textbf{わずか $k=3$ 回の反復で $\Ord{h^6}$ 以上の空間精度に到達}し，
CCD 演算子の空間精度を損なわないことを確認した．
さらに，各反復のコストが Thomas 法ベースの $\mathcal{O}(N^2)$ であるため，
Krylov 法の行列-ベクトル積コストに比して定数係数が小さい．
